\documentclass[tikz,border=6pt]{standalone}
\usepackage{amsmath,amssymb,mathtools,physics,bm,nicefrac,wasysym}
\usepackage{xcolor}

\definecolor{colone}{RGB}{180,60,60}
\definecolor{coltwo}{RGB}{60,110,180}
\definecolor{colthree}{RGB}{60,140,80}

\definecolor{accent}{HTML}{A13D2C}
\definecolor{accent2}{HTML}{2B5F6B}
\definecolor{muted}{HTML}{5A564F}
\definecolor{panel}{HTML}{F8F6F0}
\definecolor{linecol}{HTML}{D8D2C4}
\definecolor{ink}{HTML}{1E1C1A}
\definecolor{astroblue}{HTML}{1B4F72}
\definecolor{astrodark}{HTML}{17202A}
\definecolor{sunyellow}{HTML}{E8A33D}

\usetikzlibrary{calc,arrows.meta,positioning,decorations.pathreplacing,decorations.markings,angles,quotes,shapes.geometric}
\usepackage{tikz-3dplot}

\newcommand{\vecr}{\vec r}
\newcommand{\vecL}{\vec L}
\newcommand{\vecA}{\vec A}
\newcommand{\vecf}{\vec f}
\newcommand{\vecF}{\vec F}
\newcommand{\vecp}{\vec p}
\newcommand{\avg}[1]{\left\langle #1\right\rangle}
\newcommand{\vecx}{\vec x}
\newcommand{\hatn}{\hat n}
\newcommand{\rhat}{\hat r}
\newcommand{\phihat}{\hat\phi}
\newcommand{\thetahat}{\hat\theta}
\newcommand{\note}[1]{\textcolor{blue!60!black}{\textit{[#1]}}}

% Astronomical acronyms
\newcommand{\LST}{\mathrm{LST}}
\newcommand{\GST}{\mathrm{GST}}
\newcommand{\GMST}{\mathrm{GMST}}
\newcommand{\GAST}{\mathrm{GAST}}
\newcommand{\LMST}{\mathrm{LMST}}
\newcommand{\LAST}{\mathrm{LAST}}
\newcommand{\UT}{\mathrm{UT}}
\newcommand{\UTC}{\mathrm{UTC}}
\newcommand{\TAI}{\mathrm{TAI}}
\newcommand{\TT}{\mathrm{TT}}
\newcommand{\TDB}{\mathrm{TDB}}
\newcommand{\JD}{\mathrm{JD}}
\newcommand{\MJD}{\mathrm{MJD}}
\newcommand{\EoT}{\mathrm{EoT}}
\newcommand{\SNR}{\mathrm{SNR}}
\newcommand{\RON}{\mathrm{RON}}

% Helper commands for specific diagrams
\newcommand{\midlinelabel}[3]{
    \node (midlabel) at ($ (#1)!.5!(#2) $) {#3};
    \draw[stealth-] (#1) --  (midlabel);
    \draw[-stealth] (midlabel) -- (#2);
}
\newcommand{\midlinelabell}[3]{
    \node[fill = white, inner sep = 0.2pt, outer sep=3pt] (midlabel) at ($ (#1)!.5!(#2) $) {#3};
    \draw[stealth-] (#1) --  (midlabel);
    \draw[-stealth] (midlabel) -- (#2);
}
\newcommand{\midlinelabelll}[3]{
    \node[fill = white, inner sep = 0.2pt, outer sep=3pt,scale=0.6] (midlabel) at ($ (#1)!.5!(#2) $) {#3};
    \draw[stealth-] (#1) --  (midlabel);
    \draw[-stealth] (midlabel) -- (#2);
}



\begin{document}
\begin{tikzpicture}[>=Stealth, scale=1.0, font=\normalsize]
			\def\aone{1.6}
			\def\atwo{2.4}
			\coordinate (COM) at (0,0);

			% Barycentric orbits
			\draw[muted, dashed, line width=0.6pt] (COM) circle (\aone);
			\draw[muted, dashed, line width=0.6pt] (COM) circle (\atwo);
			\draw[ink!50, thin] (COM) -- (0,\aone) node[midway, right=1pt] {$a_1$};
			\draw[ink!50, thin] (COM) -- (0,-\atwo) node[midway, right=1pt] {$a_2$};
			\fill[ink] (COM) circle (1.6pt);
			\node[ink, below=2pt] at (COM) {COM};

			% Stars shown at quadrature, the orbital phase of extremal radial
			% velocity. Each star carries only its mass label, tied to its
			% own velocity arrow: above the arrow for Star 1 (top), below
			% the arrow for Star 2 (bottom).
			\fill[astroblue] (0,\aone) circle (3.2pt);
			\draw[->, astroblue, thick] (0,\aone) -- (-1.3,\aone);
			\node[astroblue, above=2pt] at (-0.65,\aone) {$m_1$};

			\fill[accent] (0,-\atwo) circle (3.2pt);
			\draw[->, accent, thick] (0,-\atwo) -- (1.3,-\atwo);
			\node[accent, below=2pt] at (0.65,-\atwo) {$m_2$};

			% Line of sight: a dashed ray to a small eye icon standing in for the
			% observer, placed level with the COM on the horizontal (line-of-sight)
			% axis, far to the left.
			\coordinate (Eye) at (-10,0);
			\draw[ink!60, dashed] (-\atwo,0) -- (Eye);
			\draw[ink!80, thick] (Eye) ++(-0.28,0) to[bend left=35] ++(0.56,0) to[bend left=35] ++(-0.56,0);
			\fill[ink!80] (Eye) circle (0.05);
		\end{tikzpicture}
\end{document}
